\section{Problems for Stage VI}
\label{problems:stage-VI}

\begin{enumerate}[label=\textbf{VI.\arabic*},leftmargin=*]
  \item Derive Eqs.~\eqref{eq:cdcc-fragment-target-coordinates} and the two
  reduced masses for a two-cluster projectile.  Check the kinetic-energy
  separation.

  \item Starting from Eq.~\eqref{eq:three-body-reaction-H}, project onto
  Eq.~\eqref{eq:cdcc-channel-function} and derive every term of
  Eq.~\eqref{eq:cdcc-coupled-radial}.

  \item Normalize a momentum-bin state.  Show strong convergence of a refined
  step-function projection on a fixed $L^2$ wave packet and explain why the
  projections do not converge in operator norm.

  \item Diagonalize a short geometric Gaussian basis for a one-dimensional
  short-range potential.  Track a bound state, a stabilized resonance-like
  pseudostate, and nonresonant pseudostates as the largest range changes.

  \item Derive the Coulomb open-channel boundary condition and the Whittaker
  closed-channel condition.  Verify the flux factor in
  Eq.~\eqref{eq:cdcc-open-boundary}.

  \item For a two-channel toy model with one closed channel, eliminate the
  closed component and compute the induced real polarization potential below
  threshold.  Continue above threshold and identify its imaginary part.

  \item Derive the complex-scaled smoothing formula
  Eq.~\eqref{eq:complex-scaled-breakup-response}.  Show explicitly which
  $\theta$ dependences cancel in a complete basis.

  \item Draw the three Jacobi sets for $\alpha+p+n$.  Construct the exchange
  symmetry under $p\leftrightarrow n$ in the isospin limit and identify which
  electromagnetic terms break it.

  \item In the $p+p+{}^7$Be model, construct symmetrized amplitudes for the two
  $p+{}^8$B arrangements.  Give a numerical test for double counting in the
  inclusive sum.

  \item Expand the spherical Bessel transforms at small recoil momentum and
  show why an $s$ wave can dominate a selected zero-recoil bin over a larger
  integrated $p$-wave component.

  \item Write the three $\alpha+n+n$ Jacobi transformations, including their
  reduced masses.  Construct basis combinations with the correct neutron
  exchange sign for $S=0$ and $S=1$.

  \item Starting from the extended complex-scaled resolution, derive
  Eq.~\eqref{eq:he6-csm-energy-response}.  Explain why the selected pole term
  can have either sign while the complete physical response is nonnegative.

  \item Represent the Riesz projector
  Eq.~\eqref{eq:he6-resonance-projector} in a real pseudostate basis and
  prove Eq.~\eqref{eq:he6-resonance-projector-pseudostates}.  Give numerical
  tests for convergence of its rank and idempotency.

  \item Derive the phase-space Jacobian in
  Eq.~\eqref{eq:he6-ddbux-definition} for one declared Jacobi momentum
  normalization.  Show that integration over both subsystem energies recovers
  the corresponding three-body breakup strength.

  \item For a two-channel incoming block, choose explicit left and right null
  vectors and verify Eq.~\eqref{eq:connection-inverse-pole}.  Identify the
  combinations of CDCC source and CSLS detector overlaps that multiply its
  residue.
\end{enumerate}
